\section{Percepción Vehicular y Reconocimiento de Matrículas}

La automatización avanzada de un parqueadero inteligente requiere trascender la
mera detección de presencia vehicular para implementar una capa de percepción
integral. En este nivel, el sistema analiza múltiples dimensiones
atributo-espaciales de los automóviles que aproximan a las barreras de acceso,
combinando la clasificación cromática y tipológica del vehículo con la
localización y lectura óptica de su placa de matrícula. Esta verificación
multi-modal no solo agiliza el flujo de paso, sino que fortalece la seguridad al
correlacionar las características físicas de la carrocería con el registro de la
matrícula almacenado en la base de datos del recinto.

\subsection{Clasificación de atributos visuales en espacios de características
discriminativos}

Como analizan Lima y Laroca \cite{lima2024vehiclecolor}, la categorización fina
de vehículos (tipo de carrocería y color predominante) proporciona una capa de
seguridad crítica en el control de acceso a parqueaderos. La identificación del
tipo de vehículo (sedán, camioneta, minibús, motocicleta) y su color permite
establecer mecanismos de auditoría cruzada: si un automóvil intenta ingresar
utilizando una placa registrada a nombre de un vehículo de distinta categoría o
tonalidad, el sistema detecta de inmediato la discrepancia y emite una alerta
por posible clonación o intercambio de matrícula.

Desde la perspectiva del procesamiento visual, la extracción de estos atributos
enfrenta desafíos ópticos significativos en el punto de acceso, como sombras
proyectadas por las casetas, reflejos solares en la pintura y variaciones en las
luminarias nocturnas. Para resolver estas ambigüedades, las arquitecturas
profundas modernas (como EfficientNet, ResNet o Vision Transformers) proyectan
los recortes del vehículo hacia espacios de características latentes de alta
dimensión. En estos espacios, la red aprende a aislar la cromaticidad intrínseca
de la pintura frente a los brillos especulares, permitiendo separar con alta
precisión clases no lineales en el espacio cromático y morfológico.

\subsection{Localización de objetos pequeños con geometría específica: detección
de matrículas}

En el encuadre general de una cámara de supervisión instalada en la barrera de
un parqueadero, la placa de matrícula ocupa únicamente una fracción diminuta del
total de píxeles del fotograma (configurando el problema clásico de detección de
objetos pequeños o \textit{small targets}). Como demuestran Laroca et al.
\cite{laroca2018robust} y Auza Banegas et al. \cite{auzabanegas2026blpr},
localizar con precisión esta subregión rectangular es una etapa crítica, pues
cualquier desviación geométrica en el recubrimiento de la placa degrada
drásticamente la tasa de éxito del módulo subsecuente de reconocimiento óptico
de caracteres (OCR).

Para superar la atenuación espacial que sufren los objetos pequeños al atravesar
las capas profundas de submuestreo, los detectores modernos apoyan su
arquitectura en Redes de Pirámide de Características (\textit{Feature Pyramid
Networks}, FPN) y en el aprovechamiento de patrones geométricos a priori. La
matrícula vehicular posee prioris estructurales altamente discriminativos, como
una relación de aspecto rectangular definida (aproximadamente $3:1$ o $2:1$),
bordes de alto contraste y una densidad interna de gradientes generada por los
caracteres alfanuméricos. Los cabezales de detección adaptativos utilizan estas
propiedades para predecir cajas delimitadoras precisas sobre la placa, aislando
la Región de Interés (ROI) de forma resistente a las variaciones de distancia
focal y altura de montaje de las cámaras.

\subsection{Del OCR clásico por segmentación al modelado de secuencias}

Como detallan las revisiones sistemáticas de Anagnostopoulos et al.
\cite{anagnostopoulos2008lpr} y Shashirangana et al. \cite{shashirangana2021alpr},
los sistemas tradicionales de reconocimiento de matrículas (ALPR) dependían de
un enfoque secuencial en tres etapas: localización de la placa, 
binarización/segmentación individual de cada carácter y clasificación mediante
coincidencia de plantillas (\textit{template matching}) o redes densas simples. 

Sin embargo, estos métodos clásicos por segmentación colapsan frecuentemente en
las condiciones operativas reales de un parqueadero al aire libre o
semiconstruido:
\begin{itemize}
    \item \textbf{Fallos por fusión de caracteres:} La presencia de suciedad,
        barro o sombras sobre la placa provoca que dos letras contiguas se unan
        en una sola máscara de binarización, imposibilitando su separación.
        
    \item \textbf{Fraccionamiento por imperfecciones:} Tornillos de sujeción,
        marcos metálicos o peladuras en la pintura generan falsos fragmentos que
        el algoritmo interpreta erróneamente como caracteres individuales.
        
    \item \textbf{Sensibilidad al umbral de binarización:} Cambios bruscos de
        iluminación entre el día y la noche impiden definir un umbral estático
        de Otsu o adaptativo que aísle limpiamente el texto sin deformarlo.
\end{itemize}

Para erradicar esta fragilidad, la visión por computadora moderna migró hacia
modelos de reconocimiento de secuencias libres de segmentación
(\textit{segmentation-free}). Arquitecturas como CRNN, propuestas por Shi, Bai y
Yao \cite{shi2017crnn}, y los modelos desarrollados por Li y Wang 
\cite{li2018reading}, tratan la lectura de la placa como un problema de
traducción de secuencia a secuencia visual, procesando la imagen de la matrícula
completa en una única pasada \textit{end-to-end} sin necesidad de cortar letra
por letra.

\subsection{Modelos de secuencias (RNN, LSTM) y alineación mediante CTC}

Como establecen Li y Wang \cite{li2018reading}, la arquitectura de un
reconocedor de matrículas libre de segmentación se compone de tres etapas
acopladas: un extractor convolucional (CNN), un modelador contextual recurrente
(RNN/LSTM) y una capa de decodificación y transcripción. El extractor CNN
transforma el recorte rectangular de la placa en una secuencia horizontal de
vectores de características alineados de izquierda a derecha. Posteriormente,
redes recurrentes bidireccionales con memoria de largo-corto plazo
(\textit{Bidirectional Long Short-Term Memory}, BLSTM) \cite{hochreiter1997lstm}
analizan la secuencia para capturar dependencias contextuales entre los
caracteres adyacentes.

El desafío matemático central en este esquema radica en que la longitud de la
secuencia de características generada por la CNN no coincide en dimensión con el
número de caracteres reales impresos en la placa. Como resolvieron Graves et al.
\cite{graves2006ctc}, el algoritmo de Clasificación Temporal Conexionista
(\textit{Connectionist Temporal Classification}, CTC) resuelve intuitivamente
este problema de alineación no supervisada sin requerir anotaciones a nivel de
píxel para cada letra.

Para entender la lógica del algoritmo CTC de forma intuitiva:
\begin{enumerate}
    \item A medida que el modelo desliza su ventana de análisis a lo largo de la
        imagen de la placa, la red emite probabilidades de caracteres para
        múltiples divisiones verticales delgadas.

    \item Como el ancho de un carácter abarca varias divisiones verticales, la
        red genera predicciones repetidas y espacios vacíos. Por ejemplo, al
        leer la secuencia ``AB1'', la salida cruda fotograma a fotograma podría
        resultar en la cadena emitidamente sobredimensionada:
        \texttt{``A-A--B---1-1''}, donde el símbolo \texttt{``-''} representa una
        etiqueta especial de espacio en blanco (\textit{blank}).

    \item El algoritmo CTC aplica una función de mapeo de colapso condicional
        ($\mathcal{B}$) que opera en dos pasos continuos: primero condensa las
        repeticiones consecutivas de caracteres idénticos no separados por un
        espacio (\texttt{``A-A--B---1-1''} $\rightarrow$ \texttt{``A-A-B-1-1''}), y
        posteriormente elimina todos los símbolos vacíos (\texttt{``-''},
        \texttt{``blank''}).
\end{enumerate}

Como resultado de esta condensación, la secuencia se transforma limpiamente en
la cadena de texto final \texttt{``AB1''}. Esta operación permite entrenar el
modelo de OCR directamente sobre imágenes de matrículas completas etiquetadas
únicamente con su texto de referencia, dotando al sistema de una robustez
extrema ante ligeras deformaciones geométricas o imperfecciones en el fondo de
la placa.

\subsection{Dominio de aplicación: estructura y gramática formal de las
matrículas en Bolivia}

Como documentan Auza Banegas et al. \cite{auzabanegas2026blpr} en el desarrollo
del marco BLPR y la normativa del Registro Único para la Administración
Tributaria Municipal (RUAT), el estándar de matrículas vehiculares en Bolivia
(vigente desde 1997) sigue una estructura alfanumérica formalmente estandarizada.
La placa convencional para automóviles se compone de tres a cuatro dígitos
numéricos seguidos de tres letras mayúsculas (formatos habituales
\texttt{``1234ABC''} o \texttt{``456DEF''}), estampadas en color azul sobre un fondo
blanco reflectante con la leyenda \texttt{``BOLIVIA''} en la parte superior.

Conocer a priori esta sintaxis matemática y gramatical permite implementar una
capa de post-procesamiento determinista basada en reglas que corrige
automáticamente posibles ambigüedades producidas por el reconocedor visual OCR:

\begin{equation}
    \text{Regla de Corrección Sintáctica} =
    \begin{cases} 
        \begin{array}{l}
            \text{'O'} \rightarrow \text{'0'}, \; \text{'I'} \rightarrow \text{'1'}, \; \text{'B'} \rightarrow \text{'8'} \\
            \text{si la posición } p \in \{1, 2, 3, 4\}
        \end{array} \\[3ex]
        \begin{array}{l}
            \text{'0'} \rightarrow \text{'O'}, \; \text{'1'} \rightarrow \text{'I'}, \; \text{'8'} \rightarrow \text{'B'} \\
            \text{si la posición } p \in \{5, 6, 7\}
        \end{array}
    \end{cases}
\end{equation}

Por ejemplo, si la cámara captura un vehículo ingresando al parqueadero y el
modelo de OCR predice la cadena ambigua \texttt{``12O4ABC''}, el módulo de
post-procesamiento evalúa la gramática oficial de Bolivia. Al detectar que la
tercera posición corresponde estrictamente al bloque numérico inicial, el
sistema fuerza la conversión del carácter alfabético ambiguo \texttt{``O''} hacia
el dígito numérico equivalente \texttt{``0''}, corrigiendo la lectura a
\texttt{``1204ABC''}. De forma análoga, si el reconocedor detecta un número
\texttt{``1''} en el bloque posterior reservado para letras, el filtro lo mapea al
carácter alfabético \texttt{``I''}. Esta validación gramatical previa a la
consulta en la base de datos incrementa significativamente la tasa de
coincidencia exacta (\textit{Exact Match}) y reduce el índice de error por
carácter (CER) en el recinto.

\subsection{Factores de degradación visual en el reconocimiento de matrículas}

En la operación continua de un parqueadero real, la calidad de la imagen
capturada en las barreras se ve afectada por múltiples factores ambientales y
mecánicos. De acuerdo con las categorizaciones de anomalías desarrolladas por
Auza Banegas et al. \cite{auzabanegas2026blpr} y Rio-Alvarez et al.
\cite{rioalvarez2019effects}, los principales factores de degradación visual se
estructuran en:

\begin{enumerate}
    \item \textbf{Distorsión por perspectiva y ángulo agudo (\textit{Skewed
        Plates}):} Representa la degradación más frecuente. Debido a que las
        cámaras suelen montarse en postes laterales o esquinas superiores de las
        casetas, la imagen de la placa sufre una deformación perspectiva
        pronunciada, deformando la geometría de los caracteres alfanuméricos.

    \item \textbf{Desenfoque por movimiento (\textit{Motion Blur}):} Ocurre
        cuando el vehículo aproxima a la barrera a una velocidad superior a la
        tasa de refresco del obturador de la cámara, dispersando los bordes de
        los caracteres a lo largo de la trayectoria.
        
    \item \textbf{Variaciones extremas de iluminación y brillos especulares
        (\textit{Low-Illumination / Glare}):} Incluye desde la escasez de luz en
        horarios nocturnos hasta el deslumbramiento provocado por las luces
        delanteras del automóvil o el sol incidente de frente, eliminando el
        contraste entre el texto azul y el fondo blanco.

    \item \textbf{Oclusiones parciales, suciedad y deterioro físico
        (\textit{Damaged / Occluded Plates}):} Placas dobladas, cubiertas por
        barro en temporadas de lluvia o parcialmente tapadas por enganches de
        remolque o marcos protectores de plástico.
\end{enumerate}

Para neutralizar estos factores sin sobrecargar la latencia del sistema, las
soluciones modernas integran etapas de preprocesamiento condicional —como
rectificación geométrica por homografía mediante detección de esquinas, ajuste
dinámico de contraste (CLAHE/Gamma) y filtros de respuesta estocástica— antes de
pasar la imagen normalizada al motor de OCR.
